% Lösungen Einheit 4 von 4 – Gleichungen aufstellen (zusammenbau v0.1)
\einheitenkopf{Einheit 4 von 4 · Gleichungen aufstellen}

%% TODO Lösung der Erklärzeile 19f) fehlt in der Bank
\erg{19}{a) $x + 9 = 23$ \quad b) $4x = 36$ \quad c) $x - 7 = 15$ \quad d) $x : 2 = 13$ \quad e) $2x + 5 = 31$ \quad f) \ldots \quad g) $x + 18 = 45$, $x = 27$ \quad h) $3x + 8 = 44$, $x = 12$ \quad i) $x + (x + 4) = 28$, $x = 12$: Der Bruder ist 12, Mia 16 Jahre alt. \quad j) $46 = 2 \cdot (8 + b)$, $23 = 8 + b$, $b = 15$ cm \quad k) $5 \cdot (380 + x) = 2\,400$; $380 + x = 480$, $x = 100$}
\erg{20}{a) $b = A : a$}
\erg{21}{a) Von der Zahl wird 6 abgezogen, nicht umgekehrt: $x - 6 = 10$, $x = 16$. \quad b) Die ganze Summe wird vervierfacht, also mit Klammer: $4(x + 3) = 36$, $x = 6$. \quad c) 4 Busse: Einen halben Bus gibt es nicht, und mit 3 Bussen fehlen Plätze – also aufrunden. \quad d) Zum Beispiel: „Ich denke mir eine Zahl und addiere 12, dann erhalte ich 30." Die Zahl ist $x = 18$. \quad e) $4 + 2{,}20x = 37$, $2{,}20x = 33$, $x = 15$: 15 km}
